%% ============================================================
%% sections/s360e-4.tex
%% United States Code, Title 21 - § 360e-4. Predetermined change control plans for devices
%% (FD&C Act § 515C). Source identifier: /us/usc/t21/s360e–4
%% Relevance to the Physical AI oncology trial context:
%% Predetermined change control plans for devices; the single best statutory anchor for an automated verification-before-change rule on AI/ML device software.
%% Reproduced faithfully from the OLRC USLM XML (through Pub. L. 119-93).
%% No bracketed instructions or file names are inserted into the law.
%% ============================================================
\uscsection{§~360e-4.}{Predetermined change control plans for devices}
\uscprov{1}{\textbf{(a)}\textbf{ Approved devices}}
\uscprov{2}{\textbf{(1)}\textbf{ In general} Notwithstanding section 360e(d)(5)(A) of this title, a supplemental application shall not be required for a change to a device approved under section 360e of this title, if such change is consistent with a predetermined change control plan that is approved pursuant to paragraph (2).}
\uscprov{2}{\textbf{(2)}\textbf{ Predetermined change control plan} The Secretary may approve a predetermined change control plan submitted in an application, including a supplemental application, under section 360e of this title that describes planned changes that may be made to the device (and that would otherwise require a supplemental application under section 360e of this title), if the device remains safe and effective without any change.}
\uscprov{2}{\textbf{(3)}\textbf{ Scope} The Secretary may require that a change control plan include labeling required for safe and effective use of the device as such device changes pursuant to such plan, notification requirements if the device does not function as intended pursuant to such plan, and performance requirements for changes made under the plan.}
\uscprov{1}{\textbf{(b)}\textbf{ Cleared devices}}
\uscprov{2}{\textbf{(1)}\textbf{ In general} Notwithstanding section 360(k) of this title, a premarket notification shall not be required for a change to a device cleared under section 360(k) of this title, if such change is consistent with an established predetermined change control plan granted pursuant to paragraph (2).}
\uscprov{2}{\textbf{(2)}\textbf{ Predetermined change control plan} The Secretary may clear a predetermined change control plan submitted in a notification submitted under section 360(k) of this title that describes planned changes that may be made to the device (and that would otherwise require a new notification), if-}
\uscprov{3}{\textbf{(A)} the device remains safe and effective without any such change; and}
\uscprov{3}{\textbf{(B)} the device would remain substantially equivalent to the predicate.}
\uscprov{2}{\textbf{(3)}\textbf{ Scope} The Secretary may require that a change control plan include labeling required for safe and effective use of the device as such device changes pursuant to such plan, notification requirements if the device does not function as intended pursuant to such plan, and performance requirements for changes made under the plan.}
\uscprov{1}{\textbf{(c)}\textbf{ Predicate devices} In making a determination of substantial equivalence pursuant to section 360c(i) of this title, the Secretary shall not compare a device to changed versions of a device implemented in accordance with an established predetermined change control plan as a predicate device. Only the version of the device cleared or approved, prior to changes made under the predetermined change control plan, may be used by a sponsor as a predicate device.}
\uscsource{(June 25, 1938, ch. 675, §~515C, as added Pub. L. 117-328, div. FF, title III, §~3308(a), Dec. 29, 2022, 136 Stat. 5835.)}
\usccrosshead{Editorial Notes}
\uscnotehead{Prior Provisions}
\uscnotepar{A prior section 515C of act June 25, 1938, was renumbered section 515B and is classified to section 360e-3 of this title.}

%% - DRAFTING INSTRUCTIONS (added for VVUQ-04 draft-bill; the reproduced statutory text above is unchanged) -
\begin{draftbox}{§ 360e-4 (21 U.S.C. § 360e-4) - predetermined change control plans; THE KEYSTONE}
\dstep{Role in the amendment: § 360e-4 is the predetermined change control plan authority, the single closest existing analogue to "validate the change before deploy." Today it is voluntary, submission-based, and not automated. This is the keystone conforming amendment: it makes an automated verification-before-change protocol mandatory for a plan that governs a device which generates or executes robot-patient interaction code, binding the plan to new section 515D. The new section 515D is inserted in the Act immediately after this section.}
\dstep{Amendatory action: Section 515C of the Federal Food, Drug, and Cosmetic Act (21 U.S.C. § 360e-4) is amended (1) in subsection (a)(3) (Scope) and in subsection (b)(3) (Scope), by inserting before the period the following: ", and, for a device that generates or executes robot-patient interaction code, an automated verification-before-change protocol consistent with section 515D that must clear before any change is implemented"; and (2) by adding at the end the following new subsection: "(d) Physical AI device software functions. A predetermined change control plan for a device that generates or executes robot-patient interaction code shall require that each planned change clear the verification process required by section 515D, in advance of generation or execution, as a condition of being implemented without a new submission under section 360e or section 360(k)."}
\dstep{Substance: the automated verification-before-change rule must carry the gate discipline from papers/VVUQ-03/final-paper/sections/statutory\_text.tex (verification fraction equal to 1.0 as an equality; validation agreement at or above the per-gate threshold; coefficient-of-variation bound; hard catastrophe predicate; the ACCEPT, BLOCK, ESCALATE decision rule), so a change that does not clear is blocked ex ante. Keep subsection (c) (predicate devices) consistent with the substantial-equivalence alignment in sections/s360c.tex.}
\dstep{Source synthesis: process papers/VVUQ-04/instruct-bill/02-fda-ai-device-regulation.md (the December 2024 PCCP final guidance and the January 2025 AI-enabled device software lifecycle draft, and the central finding that the PCCP regime is voluntary and submission-based, not mandatory, ex ante, and automated) and 01-federal-statutory-framework.md (§ 360e-4 added by the Food and Drug Omnibus Reform Act of 2022) together with papers/VVUQ-03/final-paper/sections/statutory\_text.tex and sections/policy\_memo.tex (the policy case that the predetermined change control plan is the closest analogue the bill makes mandatory). Resolve to references.bib keys usc-360e4, pl-fdora, fda-pccp-2024, and asmevv40 (ASME V\&V 40-2018 as the credibility basis).}
\dstep{Comparative print rendering: reproduce subsections (a), (b), and (c) unchanged except for the two scope insertions, which are wrapped in \cprinsert{...} inside (a)(3) and (b)(3); show the new subsection (d) wrapped in \cprinsert{...} at the end. This is the section whose marked changes most directly carry the bill's thesis, so keep the insertions precise and self-contained.}
\dstep{Formatting: single hyphens only; § for every codified reference; even RaggedRight spacing; reword so no inserted clause ends in a one- or two-word line; do not alter the reproduced subsection (c) or the source credit and prior-provisions note.}
\end{draftbox}
